\documentclass[12pt,a4paper]{article}
\usepackage{amsmath, amssymb, amsthm}
\usepackage[utf8]{inputenc}
\usepackage[T1]{fontenc}
\usepackage{geometry}
\usepackage{enumitem}
\usepackage{xcolor}
\geometry{margin=2.5cm}

\newcommand{\easy}{{\color{green!60!black}$\bigstar$}}
\newcommand{\medium}{{\color{orange}$\bigstar\bigstar$}}
\newcommand{\hard}{{\color{red}$\bigstar\bigstar\bigstar$}}
\newcommand{\pd}[2]{\dfrac{\partial #1}{\partial #2}}

\title{\textbf{Solutions: Additional Exercises 4}\\
\large Gradient, Directional Derivatives \& Higher-Order Partial Derivatives}
\date{}

\begin{document}
\maketitle
\thispagestyle{empty}

%----------------------------------------------------------
\section*{Exercise 1: Directional derivatives}
%----------------------------------------------------------

\begin{enumerate}[leftmargin=*, label=(\alph*)]

\item $f=x^2+xy-y^2$. $\nabla f=(2x+y,\,x-2y)$. At $(1,2)$: $\nabla f=(4,-3)$.
Unit vector: $\hat{u}=(3,4)/5=(3/5,4/5)$.
$D_{\hat{u}}f=(4)(3/5)+(-3)(4/5)=12/5-12/5=0$.

\item $f=xy+yz+zx$. $\nabla f=(y+z,\,x+z,\,y+x)$. At $(2,1,3)$: $\nabla f=(4,5,3)$.
$|\bar{a}|=\sqrt{9+16+144}=13$. $\hat{a}=(3,4,12)/13$.
$D_{\hat{a}}f=(4\cdot 3+5\cdot 4+3\cdot 12)/13=(12+20+36)/13=68/13$.

\item $f=e^x\cos y$. $\nabla f=(e^x\cos y,\,-e^x\sin y)$. At $(0,0)$: $\nabla f=(1,0)$.
$\hat{u}=(\cos(\pi/4),\sin(\pi/4))=(1/\sqrt{2},1/\sqrt{2})$.
$D_{\hat{u}}f=1/\sqrt{2}$.

\item $f=\ln(x^2+y^2+z^2)$. $\nabla f=\frac{2(x,y,z)}{x^2+y^2+z^2}$. At $(1,2,1)$: $r^2=6$, $\nabla f=(2/6,4/6,2/6)=(1/3,2/3,1/3)$.
$|\bar{a}|=\sqrt{4+16+16}=6$. $\hat{a}=(2,4,4)/6=(1/3,2/3,2/3)$.
$D_{\hat{a}}f=(1/3)(1/3)+(2/3)(2/3)+(1/3)(2/3)=1/9+4/9+2/9=7/9$.

\item $f=xyz$. $\nabla f=(yz,xz,xy)$. At $(2,4,2)$: $\nabla f=(8,4,8)$.
$|\bar{u}|=\sqrt{16+4+16}=6$. $\hat{u}=(4,2,-4)/6=(2/3,1/3,-2/3)$.
$D_{\hat{u}}f=8(2/3)+4(1/3)+8(-2/3)=16/3+4/3-16/3=4/3$.

\item $f=x^2y$. $\nabla f=(2xy,x^2)$. At $(1,1)$: $\nabla f=(2,1)$.
$\hat{u}=(\cos(\pi/3),\sin(\pi/3))=(1/2,\sqrt{3}/2)$.
$D_{\hat{u}}f=2(1/2)+1(\sqrt{3}/2)=1+\sqrt{3}/2$.

\item $f(x,y)=x^4y^2/(x^4+y^2)$ at $(0,0)$, $\bar{u}=(1,1)$, $\hat{u}=(1/\sqrt{2},1/\sqrt{2})$.
$D_{\hat{u}}f(0,0)=\lim_{t\to 0}\frac{f(t/\sqrt{2},t/\sqrt{2})-0}{t}$.
$f(t/\sqrt{2},t/\sqrt{2})=\frac{(t/\sqrt{2})^4(t/\sqrt{2})^2}{(t/\sqrt{2})^4+(t/\sqrt{2})^2}=\frac{t^6/8}{t^4/4+t^2/2}=\frac{t^6/8}{t^2(t^2/4+1/2)}=\frac{t^4/8}{t^2/4+1/2}$.
As $t\to 0$: $\frac{0}{1/2}=0$. So $D_{\hat{u}}f(0,0)=\lim_{t\to 0}\frac{0}{t}=0$.
\textbf{Answer: $0$.}

\end{enumerate}

%----------------------------------------------------------
\section*{Exercise 2: Finding directions}
%----------------------------------------------------------

\begin{enumerate}[leftmargin=*, label=(\alph*)]

\item $f=3x+4y$. $\nabla f=(3,4)$ (constant). Direction of steepest ascent: $\hat{u}=(3,4)/5$.
Rate (maximum directional derivative): $|\nabla f|=\sqrt{9+16}=5$.

\item $f=x^2y+y^3$. $\nabla f=(2xy,x^2+3y^2)$. At $(1,-1)$: $\nabla f=(-2,4)$.
Want $D_{\hat{u}}f=\nabla f\cdot\hat{u}=0$: $\hat{u}\perp(-2,4)$, so $\hat{u}=\pm(4,2)/\sqrt{20}=\pm(2,1)/\sqrt{5}$.

\item $f=xy+z$. $\nabla f=(y,x,1)$. At $(0,1,-2)$: $\nabla f=(1,0,1)$, $|\nabla f|=\sqrt{2}$.
$D_{\hat{u}}f=\nabla f\cdot\hat{u}$. Maximum is $|\nabla f|=\sqrt{2}<1$? No: $\sqrt{2}\approx 1.41>1$.
$D_{\hat{u}}f=1$ is achievable: need $\nabla f\cdot\hat{u}=1$, i.e.\ $\hat{u}_1+\hat{u}_3=1$.
Take $\hat{u}=(1/\sqrt{2},0,1/\sqrt{2})$... check: $(1)(1/\sqrt{2})+(0)(0)+(1)(1/\sqrt{2})=2/\sqrt{2}=\sqrt{2}\ne 1$.
Hmm. Solve: $(1,0,1)\cdot\hat{u}=1$, $|\hat{u}|=1$. Let $\hat{u}=(a,b,c)$: $a+c=1$, $a^2+b^2+c^2=1$.
One solution: $b=0$, $c=1-a$, $a^2+(1-a)^2=1\Rightarrow 2a^2-2a=0\Rightarrow a=0$ or $a=1$.
So $\hat{u}=(0,0,1)$ or $\hat{u}=(1,0,0)$.

For $D_{\hat{u}}f=20$: need $\nabla f\cdot\hat{u}=20$, but $|\nabla f\cdot\hat{u}|\le|\nabla f|=\sqrt{2}<20$. \textbf{No such direction exists.}

\item $f=x^2y+y^3$. $|\nabla f(1,-1)|=|(-2,4)|=\sqrt{20}=2\sqrt{5}\approx 4.47<20$.
\textbf{No such direction exists}, since $|D_{\hat{u}}f|\le|\nabla f|=2\sqrt{5}<20$.

\end{enumerate}

%----------------------------------------------------------
\section*{Exercise 3: Gradient at a point}
%----------------------------------------------------------

\begin{enumerate}[leftmargin=*, label=(\alph*)]

\item $u=x+y+2xy$. $u_x=1+2y$, $u_y=1+2x$. At $(1,1)$: $\nabla u=(3,3)$.

\item $u=x^2+y^2+z^2$. $\nabla u=(2x,2y,2z)$. At $(2,0,0)$: $\nabla u=(4,0,0)$.

\item $u=\ln(x^2+y^2+z^2)$. $\nabla u=\frac{2(x,y,z)}{x^2+y^2+z^2}$. At $(1,1,1)$: $r^2=3$, $\nabla u=(2/3,2/3,2/3)$.

\item $u=x^2e^{y-z}$. $u_x=2xe^{y-z}$, $u_y=x^2e^{y-z}$, $u_z=-x^2e^{y-z}$.
At $(2,1,1)$: $e^0=1$. $\nabla u=(4,4,-4)$.

\item $u=\arctan(y/x)+\arctan(x/y)$. Both are functions of $y/x$ and $x/y=1/(y/x)$.
$\arctan(t)+\arctan(1/t)=\pi/2$ for $t>0$, $-\pi/2$ for $t<0$ (constant on each quadrant).
So $u$ is constant (piecewise), $\nabla u=\mathbf{0}$.

\item $u=\arctan(xyz)+\operatorname{arccot}(xyz)$. Since $\arctan(t)+\operatorname{arccot}(t)=\pi/2$ for all $t$, $u=\pi/2$ is constant, so $\nabla u=\mathbf{0}$.

\item $u=f(r)$, $r=\sqrt{x^2+y^2+z^2}$. $u_x=f'(r)\cdot\frac{x}{r}$.
$\nabla u=f'(r)\cdot\frac{(x,y,z)}{r}=\frac{f'(r)}{r}\bar{r}$.

\end{enumerate}

%----------------------------------------------------------
\section*{Exercise 4: Second-order partial derivatives}
%----------------------------------------------------------

\begin{enumerate}[leftmargin=*, label=(\alph*)]

\item $z=\sin(xy)$. $z_x=y\cos(xy)$. $z_{xy}=\cos(xy)-xy\sin(xy)$.
$z_{xyy}=-x\sin(xy)\cdot 1+(-x\sin(xy)-x^2y\cos(xy))$\ldots Let's compute systematically:
$\frac{\partial^2 z}{\partial x\,\partial y^2}=\frac{\partial}{\partial x}(z_{yy})$.
$z_y=x\cos(xy)$. $z_{yy}=-x^2\sin(xy)$. $z_{x(yy)}=-2x\sin(xy)-x^2y\cos(xy)$.

\item $z=x^3y+xy^3$. $z_x=3x^2y+y^3$, $z_{xy}=3x^2+3y^2$.
$z_y=x^3+3xy^2$, $z_{yx}=3x^2+3y^2$. $f_{xy}=f_{yx}=3x^2+3y^2$. $\checkmark$

\item $z=x^y=e^{y\ln x}$. $z_x=yx^{y-1}$, $z_{xy}=x^{y-1}+y x^{y-1}\ln x=x^{y-1}(1+y\ln x)$.
$z_y=x^y\ln x$, $z_{yx}=yx^{y-1}\ln x+x^y/x=x^{y-1}(y\ln x+1)$. Equal. $\checkmark$

\item $u=\arctan(y/x)$. $u_x=-y/(x^2+y^2)$.
$u_{xx}=\frac{2xy}{(x^2+y^2)^2}$.
$u_y=x/(x^2+y^2)$.
$u_{yy}=\frac{-2xy}{(x^2+y^2)^2}$.
$u_{xx}+u_{yy}=0$. $\checkmark$

\item $u=e^x(x\cos y-y\sin y)$. $u_x=e^x(x\cos y-y\sin y)+e^x\cos y=e^x((x+1)\cos y-y\sin y)$.
$u_{xx}=e^x((x+2)\cos y-y\sin y)$.
$u_y=e^x(-x\sin y-\sin y-y\cos y)=e^x(-(x+1)\sin y-y\cos y)$.
$u_{yy}=e^x(-(x+1)\cos y-\cos y+y\sin y)=e^x(-(x+2)\cos y+y\sin y)$.
$u_{xx}+u_{yy}=e^x[(x+2)\cos y-y\sin y-(x+2)\cos y+y\sin y]=0$. $\checkmark$

\item $u=\ln(1/r)=-\ln r$, $r=\sqrt{(x-a)^2+(y-b)^2}$.
$u_x=-r_x/r=-(x-a)/r^2$. $u_{xx}=-r^2-(x-a)2r\cdot r_x)/r^4$... more carefully:
$u_{xx}=\frac{-(r^2-(x-a)^2\cdot 2)/r^2}{r^2}$\ldots

Let $X=x-a$, $Y=y-b$, $r^2=X^2+Y^2$. $u=-\frac{1}{2}\ln(X^2+Y^2)$.
$u_X=-X/(X^2+Y^2)$. $u_{XX}=\frac{-(X^2+Y^2)+2X^2}{(X^2+Y^2)^2}=\frac{Y^2-X^2}{(X^2+Y^2)^2}$.
$u_{YY}=\frac{X^2-Y^2}{(X^2+Y^2)^2}$.
$u_{XX}+u_{YY}=0$. $\checkmark$

\end{enumerate}

%----------------------------------------------------------
\section*{Exercise 5: Second-order differential}
%----------------------------------------------------------

\begin{enumerate}[leftmargin=*, label=(\alph*)]

\item $u=x^2y^3$. $u_{xx}=2y^3$, $u_{xy}=6x y^2$, $u_{yy}=6x^2y$.
$d^2u=2y^3(dx)^2+12xy^2\,dx\,dy+6x^2y(dy)^2$.

\item $u=e^{xy}$. $u_{xx}=y^2e^{xy}$, $u_{xy}=e^{xy}+xye^{xy}=(1+xy)e^{xy}$, $u_{yy}=x^2e^{xy}$.
$d^2u=y^2e^{xy}(dx)^2+2(1+xy)e^{xy}\,dx\,dy+x^2e^{xy}(dy)^2$.

\item $u=x^y$. $u_{xx}=y(y-1)x^{y-2}$, $u_{xy}=x^{y-1}(1+y\ln x)$, $u_{yy}=x^y(\ln x)^2$.
$d^2u=y(y-1)x^{y-2}(dx)^2+2x^{y-1}(1+y\ln x)\,dx\,dy+x^y(\ln x)^2(dy)^2$.

\item $u=x\ln(xy)=x\ln x+x\ln y$. $u_x=\ln x+1+\ln y=\ln(xy)+1$. $u_y=x/y$.
$u_{xx}=1/x$. $u_{xy}=1/y$. $u_{yy}=-x/y^2$.
$d^2u=\frac{(dx)^2}{x}+\frac{2\,dx\,dy}{y}-\frac{x(dy)^2}{y^2}$.

\item $u=\sin(x+y)+\cos(x-y)$. $u_x=\cos(x+y)-\sin(x-y)$. $u_y=\cos(x+y)+\sin(x-y)$.
$u_{xx}=-\sin(x+y)-\cos(x-y)$. $u_{xy}=-\sin(x+y)+\cos(x-y)$. $u_{yy}=-\sin(x+y)-\cos(x-y)$.
$d^2u=[-\sin(x+y)-\cos(x-y)](dx)^2+2[-\sin(x+y)+\cos(x-y)]\,dx\,dy+[-\sin(x+y)-\cos(x-y)](dy)^2$.

\item $u=f(t)$ where $t=x^2-y^2$. $u_x=f'(t)\cdot 2x$. $u_{xx}=f''(t)\cdot 4x^2+2f'(t)$.
$u_y=f'(t)\cdot(-2y)$. $u_{yy}=f''(t)\cdot 4y^2-2f'(t)$.
$u_{xy}=f''(t)\cdot 2x\cdot(-2y)=-4xyf''(t)$.
$d^2u=(4x^2 f''+2f')(dx)^2-8xy f''\,dx\,dy+(4y^2 f''-2f')(dy)^2$.

\item $u=\frac{1}{2}\ln(x^2+y^2)$. $u_x=x/(x^2+y^2)$, $u_{xx}=(y^2-x^2)/(x^2+y^2)^2$, $u_{yy}=(x^2-y^2)/(x^2+y^2)^2$.
$u_{xx}+u_{yy}=0$. $\checkmark$
$u_{xy}=-2xy/(x^2+y^2)^2$.
$d^2u=\frac{y^2-x^2}{(x^2+y^2)^2}(dx)^2-\frac{4xy}{(x^2+y^2)^2}\,dx\,dy+\frac{x^2-y^2}{(x^2+y^2)^2}(dy)^2$.

\end{enumerate}

\end{document}
